\section{Loop-Level Recursion Relations}
\label{sec:loop-level-recursion}

\textbf{Literature.} Extending recursion to loop level proceeds through unitarity: loop integrands are reconstructed from on-shell tree data glued across unitarity cuts.
For the basic cutting rules, which relate discontinuities or imaginary parts of loop amplitudes to products of on-shell lower-order amplitudes, see Schwartz's discussion in \autocite[\S24.1.2]{schwartzQFT2014}.
The on-shell and generalized-unitarity program for one-loop QCD amplitudes is reviewed in \autocite{bernOnShellMethods2007}, while \autocite{Abreu:2014cla} analyzes how multiple unitarity cuts encode the discontinuities of the underlying Feynman integrals. The relation between loops and trees is made systematic through the \textit{Feynman tree theorem} in \autocite{brandhuberTreesLoops2006} and through forward limits of tree amplitudes in \autocite{Caron-Huot:2010fvq}. In the planar limit these ideas culminate in the all-loop BCFW recursion for the integrand of planar $\mathcal{N}=4$ SYM \autocite{Arkani-Hamed:2010zjl}, discussed further below.

\subsection{Planar Limit}
\label{subsec:planar-limit}

See \autocite{Arkani-Hamed:2010zjl}, where loop-level BCFW recursion is established for Supersymmetric Yang-Mills theory \textit{in the planar limit} (i.e., $SU(N)$ gauge theory with $N \to \infty$). In the planar limit, all Feymann diagrams are planar; its significance is also explained in \autocite{Beisert:2010jr}.